\documentclass[11pt]{amsart}
\usepackage{geometry} % See geometry.pdf to learn the layout options. There are lots.
\geometry{letterpaper} % ... or a4paper or a5paper or ...
%\geometry{landscape} % Activate for for rotated page geometry
%\usepackage[parfill]{parskip} % Activate to begin paragraphs with an empty line rather than an indent
\usepackage{graphicx}
\usepackage{amssymb}
\usepackage{epstopdf}
\DeclareGraphicsRule{.tif}{png}{.png}{convert #1 dirname #1/basename #1 .tif`.png}

\title{Problema 4.15}
%\date{}
% Activate to display a given date or no date

\begin{document}
\maketitle
%\section{}
%\subsection{}
\noindent Una empresa automotriz necesita un programa para manejar los montos de ventas de sus N sucursales, a lo largo de los últimos M años. Los datos son dados de esta forma:

\begin{center}
M, M \\
$MOMTO_{i,i}, MOMTO_{1,2}$, . . .  , $MOMTO_{1,N}$ \\
$MOMTO2_{2,1}, MOMTO_{2,2}$, . . .  , $MOMTO_{2,N} $\\
$MOMTO_{M,1} , MOMTO_{2,2}$, . . .  , $MOMTO_{M,N}$\\
\end{center}

\noindent Donde:
\begin{itemize}
    \item M es una variable de tipo entero que representa el número de años, 1 $ \leq $ M $ \leq $ 10.
     \item N es una variable de tipo entero que expresa el número de sucursales de la empresa, 1 $ \leq $ N $ \leq $ 35.
    \item \(MONTO_{i, j}\) es una variable de tipo real que representa lo que se vendió en el año i en la sucursal j.\end{itemize}

\noindent La información que necesitan los directores de la empresa para tomar decisiones es la siguiente:

\begin{enumerate}
    \item[a)] Sucursal que más ha vendido en los M años.
    \item[b)] Promedio de ventas por año.
    \item[c)] Año con mayor promedio de ventas.
\end{enumerate}

\end{document}